\documentclass[12pt,a4paper,final,abstract=true,headings=small]{scrartcl}
\usepackage{startup}
%%%%%%%%%%%%%%%%%%%%%%%%%%%%%%%%%%%%
% \usepackage[disable,nolists, tablesfirst, heads]{endfloat}  % Format where figures and tables will go
%%%%%%%%%%%%%%%%%%%%%%%%%%%%%%%%%%%%%%%%%%%%
\begin{document}
% \maketitle
\begin{titlepage}
    \begin{center}
        \textbf{\Huge Evaluating effort as a management tool for small game.} \\
        \vspace*{2cm}
        \textbf{\Large Using distance sampling with two dynamic models, the Gompertz and Hazard.} \\
        \vspace*{5cm}
        \emph{\Large Tomas Willebrand} \\
        \today
    \end{center}
\end{titlepage}
%% =============================================




\section{Introduction}
\input{01_IntroductionV1.tex}

%%%%%%%%%%%%%%%%%%%%%%%%%%%%%%%%%%%%%%%%%%%%%%%%%%%%%%%%%%%%%%%%%%%%%%%%%%%%%%%%%%%%%%%%%

\section{Methods}

\input{02_MethodsV2.tex.tex}
%\input{Methods_AV1.tex}
% % \subsection*{Model development and Validation}
% \input{Methods_StaticV4.tex}
% % \subsection*{Simulation --- Harvest effects}
% \input{Risk_analysis_GomprtzV3_short.tex}
% \input{Methods_Simulations_shortV1.tex}

%%%%%%%%%%%%%%%%%%%%%%%%%%%%%%%%%%%%%%%%%%%%%%%%%%%%%%%%%%%%%%%%%%%%%%%%%%%%%%%%%%%%%%%%%%%


\section*{Results}

\input{03_ResultsV2.tex}

\section*{Discussion}

Can threhold be used to set harvest quotas?

The use of complementary models to evaluate the effects of harvest on population dynamics, and the potential for using such models to inform management decisions.

The need to separate adult and juvenile survival in combination with productivity in order to better understand the effects of harvest on population dynamics.

High natural mortality and stochastic fluctuations in productivity in combination with moderate density dependent feedback.

Willow ptarmigan as a model system for understanding the effects of harvest on population dynamics, and the potential for applying insights from this system to other species and contexts.

\section{Appendix}

\input{02_Appendix_Methods.tex}

\end{document}
